\section{K-026: cyclically stable components on the effective free locus}

Write $V_n=(\mathbb C^n)^{\otimes3}$, $G_n=SL_n^3$, and
$X_n=V_n/\!/G_n$, with the degree-$ns$ to degree-$s$ regrading.
Let
\[
 S=\{(aI,bI,cI):abc=1\},\qquad
 \overline G=GL_n^3/S,
 \qquad \chi(A,B,C)=\det A\det B\det C.
\]
The kernel of $\chi$ on $\overline G$ is the effective image
of $G_n$.  Say that a quotient point is \emph{effectively free}
if the closed orbit over it has stabilizer in $G_n$ equal to the
ineffective kernel $G_n\cap S$.
Let $\tau$ cyclically permute the three tensor factors, put
$M=2^{n-1}$, and let $Y_n$ be the chain curve of the supplied
finding P-F60.

\begin{proposition}[K-026a]
Suppose $n\ge4$.  If an irreducible component $Y$ of the reduced
fixed locus $X_n^{\mu_M}$ has a nonempty open subset of effectively
free points and satisfies $\tau Y=Y$, then $Y=Y_n$.
Consequently every other such component has an $S_3$ orbit of
size three or six.  In particular, an orbit of size two cannot
occur among these components.
\end{proposition}

\begin{proof}
Fix a primitive $M$th root $\xi$.  For a nonzero effectively free
point $y$, choose a closed-orbit representative $T$.  The image
$H_T$ of $\operatorname{Stab}_{GL_n^3}(T)$ in $\overline G$
injects into $\mathbb C^*$ under $\chi$: its kernel is exactly
the effective special-linear stabilizer.  It is finite, since a
nonzero positive-degree invariant at $T$ forces a positive power
of $\chi$ to be one on the stabilizer.  Thus $H_T$ is cyclic.

If $y$ is fixed by the grading element $\xi$, the uniqueness of
the closed orbit in a reductive quotient supplies an element
of $H_T$ of character $\xi$.  Injectivity makes this element
unique; call it $h_y$.  Its order is exactly $M$.  Its
$\overline G$ conjugacy class is independent of the representative
$T$.

There are finitely many conjugacy classes of elements of order
dividing $M$ in $\overline G$: a finite-order element lies in a
maximal torus, whose $M$-torsion is finite, and conjugacy is then
quotiented by a finite Weyl group.  More explicitly in this
setting, lift such an element to $(A,B,C)\in GL_n^3$.  Its
$M$th power lies in $S$.  The $M$th-power map on the torus $S$
is surjective, so multiplication by an element of $S$ makes the
lift have $M$th power equal to the identity.  It is then a
diagonal triple with eigenvalues in $\mu_M$, up to conjugacy.

The sets of effectively free fixed points with any one of these
classes are constructible: they are Chevalley images under the
quotient map of the corresponding fixed vector spaces intersected
with the stable, effectively free locus in $V_n$.  Thus only
closed-orbit representatives enter these images.  They partition
the effectively free fixed locus because
$h_y$ is unique.  On the irreducible $Y$, one class therefore
occurs on a dense open subset.  If $\tau Y=Y$, its cyclic
permutation must be the same class, since two dense open subsets
of $Y$ meet.

Take an order-$M$ diagonal lift of that class and denote the
three eigenvalue multisets, in additive notation modulo $M$,
by $P_1,P_2,P_3$.  Equality of its conjugacy class with its cyclic
permutation, in $\overline G$, implies that all three multisets
are scalar translates of a common multiset $P$; write
$P_j=P+a_j$.  Each $a_j$ lies in $\mathbb Z/M\mathbb Z$ because
the eigenvalues of the lift lie in $\mu_M$.  As $3$ is invertible
modulo $M$, choose $a$ with $3a=a_1+a_2+a_3$.  Multiplication in
factor $j$ by the scalar of exponent $a-a_j$ has product one,
and therefore belongs to $S$.  The new lift has identical
multisets in all three factors.  Conjugating independently
within each factor makes it $(A,A,A)$, where $A^M=I$ and
\[
                         (\det A)^3=\xi.
\]
In particular, $\det A$ has order $M$.

Let the diagonal exponents of $A$ be $w_1,\ldots,w_n$.
The support of any tensor fixed by $(A,A,A)$ consists of
ordered triples with $w_i+w_j+w_k=0$.  A nonzero polystable
tensor is semistable under the product of the three diagonal
determinant-one tori,
so its support has nonnegative masses with all three mode
margins equal to one.  On forgetting the ordering of a triple
and dividing these masses by three, their exponent columns
express $\mathbf1$ as a nonnegative rational combination of
cubic exponent vectors.  Thus the complete cubic fixed support
also has $\mathbf1$ in its rational span.

Apply the complete supplied cubic-support theorem P-F30,
including its rank-deficient and arbitrary-lift cases.  Since
the determinant order is the extremal $M=2^{n-1}$, the complete
cubic support, after relabeling, is exactly
\[
                    2e_i+e_{i+1}\quad(i<n),\qquad 3e_n.
\]
The complete tensor fixed support is therefore the three ordered
permutations of each $(i,i,i+1)$ and the triple $(n,n,n)$.

For completeness, balancing on this support forces all its
coefficients to occur.  Let $x_i,y_i,z_i$ be the masses on those
three ordered triples, with the repeated index occupying the
three possible pairs.  At the first vertex, the three margin
equations give $x_1=y_1=z_1=1/2$.  Inductively they give
\[
 x_i=y_i=z_i=r_i,\qquad
 r_0=0,\quad r_i=(1-r_{i-1})/2\quad(i<n),
\]
and the final diagonal mass is $1-r_{n-1}$.  These masses are
strictly positive and unique.  The same homogeneous calculation
shows that the $3n-2$ incidence columns are linearly independent.
Hence the diagonal torus of $GL_n^3$ acts transitively on choices
of nonzero coefficients: the corresponding monomial map of
complex tori has full column rank and is surjective.  Every
such tensor belongs to the chain's $GL_n^3$ orbit.

Its quotient points form the single nonzero grading orbit
whose closure is $Y_n$.  A dense open subset of $Y$ is therefore
contained in $Y_n$, so $Y=Y_n$.  P-F60 shows that this curve
indeed lies in the effectively free locus away from the origin
and is fixed by all of $S_3$.  Any subgroup of $S_3$ containing
the three-cycle has index one or two.  Those two orbit sizes
are thus excluded for every other component under consideration,
leaving sizes three and six.
\end{proof}

\paragraph{The precise remaining cancellation gate.}
This proposition does not assert that every fixed component
meets the effectively free locus.  It also does not make the
leading coefficients of the other components integral at three.
Even for a smooth fixed curve, a grading stabilizer of order
$3^aM$ can give a normalized Hilbert density with denominator
$3^a$.  Therefore a three-element orbit of components need not
have contribution divisible by three in a three-adic valuation
ring.  K-025 constructs an actual polystable odd lift at the
requested dyadic order, so a general exclusion of these lifts
would be false.  K-026a reduces the geometry but does not close
this arithmetic gate.

\begin{lemma}[K-026b: the common phase line]
Let $y$ be an effectively free point of $X_n$ fixed by a finite
grading element $\xi$.  Suppose it lies on a fixed component
that is smooth of dimension $d$ at $y$.  Put
$D=\dim X_n=n^3-3n^2+3$.  In the point-pullback convention of
P-F23, the determinant of the moving tangent representation is
$\xi^{n^2}$.  If its local leading factor is
\[
 C=\frac{e}{\prod_{\lambda\ne1}(1-\lambda)},\qquad e>0,
\]
then
\[
                   \overline C=(-1)^{D-d}\xi^{n^2}C.
                                                        \tag{K-026b}
\]
Thus contributions from smooth components of the same dimension
lie on one real phase line.  This does not select a common sign
on that line.
\end{lemma}

\begin{proof}
At an effectively free point the quotient tangent representation
is the tensor representation minus the special-linear orbit
representation.  For a fixing lift $(A,B,C)$, the determinant
on the tensor space is
$(\det A\det B\det C)^{n^2}=\xi^{n^2}$.  Conjugation on each
$\mathfrak{sl}_n$ has determinant one.  Fixed tangent eigenvalues
are one, so removal of those directions leaves the same
determinant on the $D-d$ moving directions.  Since the eigenvalues
have absolute value one,
\[
 \overline{\prod(1-\lambda)}
 =\prod(1-\lambda^{-1})
 =(-1)^{D-d}\xi^{-n^2}\prod(1-\lambda).
\]
Taking reciprocals and using the real positive $e$ proves the
formula.  The opposite signs of the actual $n=4$ components in
K-005 show why the phase-line conclusion cannot be strengthened
to a same-sign statement without additional information.
\end{proof}
